\usepackage{amssymb,amsmath,amscd, amsthm, array}
\usepackage{cases} 
\usepackage{graphicx}
\usepackage{wrapfig}
\usepackage{color}
 \usepackage{bm}
 
 %%%%%%%%%%%%%
%\usepackage{mathpazo}
\usepackage{newtxmath,newtxtext}
%\usepackage{setspace}
\usepackage{amsmath,mathtools}
\usepackage{amsfonts}
\usepackage{tikz-cd}
%%%%%%%%%%%%%

%%%%%%%%%%%%%
\usepackage{hyperref}
\hypersetup{colorlinks}
\hypersetup{
    %bookmarks=true,         % show bookmarks bar?
    colorlinks=true,       % false: boxed links; true: colored links
    linkcolor=red,          % color of internal links (change box color with linkbordercolor)
    citecolor=magenta,        % color of links to bibliography
    filecolor=violet,      % color of file links
    urlcolor=blue      % color of external links
}
%%%%%%%%%%%%%


%%%%%%%%%%%%%%%%%%%%%%
\makeatletter
\@addtoreset{equation}{section}
\def\theequation{\thesection.\arabic{equation}}
\makeatother
%%%%%%%%%%%%%%%%%%%%%%%%%%%%%%%%%%%%%%%5
\textwidth=14.5cm
\textheight=21.5cm
\oddsidemargin=0.8cm
\evensidemargin=0.8cm
\topmargin=1cm
%%%%%%%%%%%%%%%%%%%%%%%%%%%%%%%%%%%%%%
%%%%%%%%%%%%%%%%%%%%%%%%%%%%%%%%%%%%
  \makeatletter
  \parsep   = 0pt
  \labelsep = 0pt
  \def\@listi{%
  \leftmargin = 0pt \rightmargin = 0pt
     \labelwidth\leftmargin \advance\labelwidth-\labelsep
   %  \topsep     = .5\baselineskip
     \partopsep  = 0pt \itemsep       = 0pt
     \itemindent = 0pt \listparindent = 0pt}
  \let\@listI\@listi
  \@listi
  \def\@listii{%
     \leftmargin = 0pt \rightmargin = 0pt
     \labelwidth\leftmargin \advance\labelwidth-\labelsep
     \topsep     = 0pt \partopsep     = 0pt \itemsep   = 0pt
     \itemindent = 0pt \listparindent = 0pt}
  \let\@listiii\@listii
  \let\@listiv\@listii
  \let\@listv\@listii
  \let\@listvi\@listii
  \makeatother
  
  %%%%%%%%%%%%%%%%%%%%%%%%%%%%%%%%%%%%%%


\renewcommand{\arraystretch}{1.2}
\newtheorem{theorem}{Theorem}[section]
\newtheorem{corollary}{Corollary}[theorem]
\newtheorem{lemma}[theorem]{Lemma}
\newtheorem{conjectural theorem}{Conjectural theorem}[section]
\newtheorem*{definition}{Definition}
\newtheorem*{remark}{Remark}
\newtheorem*{proposition}{Proposition}
\newtheorem*{example}{Example}

%newcommand
\newcommand{\C}{{\mathbb{C}}}
\newcommand{\F}{{\mathbb{F}}}
\newcommand{\N}{{\mathbb{N}}}
\newcommand{\Q}{{\mathbb{Q}}}
\newcommand{\R}{{\mathbb{R}}}
\newcommand{\Z}{{\mathbb{Z}}}
\newcommand{\Fh}{{\mathcal{F}}}
\newcommand{\Ah}{{\mathcal{A}}}
\renewcommand{\det}{{\mathrm{det}}}
\newcommand{\id}{{\mathrm{id}}}
\newcommand{\tr}{{\mathrm{tr}}}
\newcommand{\Sp}{{\mathrm{Sp}}}
\newcommand{\Arg}{{\mathrm{Arg}}}
\newcommand{\sdim}{{\mathrm{sdim}}}
